\documentclass[a4paper,12pt]{article}
\pagestyle{empty}

\begin{document}
\begin{center}
{\Large\bf 10112 Advanced Quantum Mechanics}
\vspace{.5cm}\\
{\large\bf Hubbard model on a ring.}
\end{center}

\def\pd{{\phantom\ldagger}}
\def\bra#1{\langle\,#1\,\vert}
\def\ket#1{\vert\,#1\,\rangle}
\def\ldagger{{\dagger}}

\begin{description}

\item[The Hubbard Model for two electrons on a ring:]\footnote{This problem set was originally written by Jim Sethna at Cornell
University.}
This problem studies the Hubbard model in one dimension.
The Hubbard model approximates the long-range Coulomb repulsion between 
electrons as an on-site repulsion energy $U.$  The model is otherwise
a tight-binding Hamiltonian with nearest-neighbor hopping.

The Hamiltonian is written
\begin{equation}
  \label{eq:2}
H=H_{TB}+H_C=-\sum_{n,s}t(c^\ldagger_{n+1,s}c^\pd_{n,s}
	  	      + c^\ldagger_{n,s}c^\pd_{n+1,s}) +
\sum_{n} U c^\ldagger_{n,\uparrow}c^\pd_{n,\uparrow}
	c^\ldagger_{n,\downarrow}c^\pd_{n,\downarrow}
  \end{equation}

$n$ is running from $1$ to $L$, where we imagine the L sites organized in
a ring so that the electron can also hop from site $L$ to site $1$.
We can see that $t$, the hopping matrix element, represents a kinetic 
energy for the electrons which is lowered if the wavefunction delocalizes
into long-wavelength plane waves.  Indeed, the eigenstates in the absence
of electron-electron interactions are plane waves.

We can introduce the creation operator for a single-particle plane-wave:
\begin{equation}
\psi^\ldagger_{k,s} = {1\over\sqrt{L}}\sum_n
e^{ink}c^\ldagger_{n,s}
\end{equation}
It can be shown (you are welcome to do that if you like) that the
Hamiltonian in this representation is diagonal and equal to

$$H_{TB}=\sum_{k,s}\varepsilon_k \psi^\ldagger_{k,s}\psi_{k,s}$$

with $\varepsilon_k = - 2 t \cos(k)$ and where $k$ runs over
$\frac{2\pi}{L}p, p=0,1,\ldots,L-1$.

\begin{itemize}
\item Discuss the meaning of this expression for $H_{TB}$.

\item Explain what the operator $H_c$ does to a state.
\end{itemize}

We will study the two-electron $N=2$ problem on a ring of 
$L = 6$ sites, with one spin up and one spin down ($S_z=0$).  We can work
in a plane-wave basis: 
\begin{equation}
  \label{eq:1}
|k_1 \uparrow, k_2 \downarrow \rangle =
\psi^\dagger_{k_1\uparrow}\psi^\dagger_{k_2\downarrow} |0\rangle
\end{equation}
The total crystal momentum is a good quantum number and we shall
consider states with zero crystal momentum so $k_1 + k_2 = 0$.

\begin{itemize}
\item  Discuss that the kinetic
term $H_{TB}$ is diagonal in this basis, and show that the on-site repulsion
has matrix elements 
$$\bra{k\uparrow,-k\downarrow}H_C\ket{k'\uparrow,-k'\downarrow} = U/L$$
independent of $k$ and $k'.$
\end{itemize}
You now have to do a bit of numerical work. You can for example use
the IPython notebook ``hubbard\_on\_a\_ring.ipynb''.

\begin{itemize}
\item Compute the eigenfunctions and eigenenergies for total crystal momentum zero.
  Plot the eigenenergies vs. $U$ for $U/t$ ranging from zero to four.
\end{itemize}

Notice that the ground state and the highest excited state change
smoothly as we increase $U$.
\begin{itemize}
\item What happens to the other two states?  Which of the states are
  singlets, and which triplets (total spin zero and one)?  Why are the
  triplet energies so boring?
\end{itemize} The coulomb repulsion alters the wavefunctions
rather more than is apparent in the energy.
\begin{itemize}
\item Calculate the ground state wavefunction in real space
  $\Psi_0(n_1,n_2)$.  It should depend only on $n_1-n_2$; plot
  $\Psi_{0}(n_1-n_2)$ and discuss the hole at $n_1 = n_2.$ How does
  the depth of the hole depend on $U$? Why?
\end{itemize}
\end{description}
\end{document}


%%% Local Variables:
%%% mode: latex
%%% TeX-master: t
%%% End:
